\section{Discussion}

Two quantities organise the results. Structural headroom measures the
opportunity that state conditioning creates: it is the loss reduction available
to a sequential rule over a perfect standalone ranking, and it is a property of
the prediction task rather than of a particular router. The controlled study
shows that redundancy creates this opportunity, and the multi-target study
shows that it is positive on every dataset and every target we evaluated.
Predictor response and ranking quality then determine how much of that
opportunity a learned router can realise. Cached representations of context
content do not carry the information that the marginal utility depends on,
while a compact summary of the predictor's response does; and as candidate
pools grow, the utility gap between the best alternatives shrinks, so a model
with stable average error becomes less able to identify the best candidate.
Both effects are visible in the results and both are addressed by the design of
R-MUR.

The multi-target results also show that realisation is uneven. METR-LA converts
the least headroom into realised gain, with the lowest oracle recovery and the
fewest targets whose gain is positive, while PEMS-BAY shows the widest spread
across targets. Structural opportunity alone therefore does not determine how
much value a router obtains, which is why the two quantities are reported
separately.

R-MUR uses greedy one-step selection, so it does not capture complementarity
that emerges only from selecting a group of individually unhelpful contexts.
\subsection{Deciding whether to route at all}

Because the improvement is not uniform across benchmarks, a practitioner needs a
way to decide before committing. The measurement protocol of the companion
studies supplies one, and its operating characteristics are measured rather than
asserted: on pilots of $n$ targets drawn from six deployments, a pilot of $24$
targets reproduces the full-deployment verdict $96\%$ of the time, with $98\%$
sensitivity and $96\%$ specificity for the decision to route at all, while
sixteen targets give $84\%$ and $86\%$. The direction of the effect is much
easier to recover than its significance. The procedure is to label a pilot of
roughly two dozen targets, run the protocol unchanged on it --- screening,
budget, ranking objective, and both the router and the pool-everything baseline
--- and route only if the paired target-level interval excludes zero above.
Cheaper substitutes do not work, and each is reported as a negative in the
companion work: a fitted predictability probe is indistinguishable from chance,
structural headroom does not order the benefit, and the cached screen is
uncorrelated with the true marginal exactly when the decision is close.

A pilot is also labelled earlier than the period it is used for, and that
shift is measured rather than assumed: scoring the same 32 targets on the
calibration and test periods of six deployments gives a within-dataset
correlation of $+.596$ between the two, per-target sign agreement of $81\%$, and
dataset-level verdict agreement of four of six --- direction transfers, the
significance call less well, and no pilot recommended routing where routing
turned out to hurt.

The pilot can also be audited against itself, and the audit is worth more than
the raw verdict: splitting it in half, pilots whose halves both resolve in the
same direction agree with the full-sample verdict $99.0\%$ of the time at
sixteen targets against $78.1\%$ for the rest. The check certifies rather than
produces a verdict --- it holds for only about a third of pilots --- so the rule
is to act when it passes and to enlarge the pilot when it does not.

Training also requires counterfactual predictor evaluations to construct the
marginal-utility targets. These costs arise offline, while deployment limits
predictor evaluations through cached screening. The current experiments study
subset-capable predictors on traffic forecasting; extending the same decision
formulation to other predictor families and context-selection domains is a
natural direction for further study.
